\documentclass[11pt]{article}
\usepackage[a4paper,margin=0.7in]{geometry}
\usepackage{amsmath,amssymb}
\usepackage[hidelinks]{hyperref}
\setlength{\parindent}{0pt}
\setlength{\parskip}{0.45em}
\setlength{\emergencystretch}{2em}
\newcommand{\dd}{\,\mathrm{d}}
\begin{document}
\begin{center}
{\Large Ramp metering: the general total-time derivative}\\[0.3em]
2 October 2026 --- explanation of the current accounting identity
\end{center}

Let \(\theta\) be one parameter of the open-loop meter-rate schedule. Let
\(P(t;\theta)\) count \emph{all} traffic still in the accounting system:
mainline, ramp, post-meter storage and external waiting. Let \(A(t)\) be
fixed exogenous arrival rate and \(E(t;\theta)\) total \emph{final exit rate}.
\(E\) is neither cumulative exits nor local ramp merge entry. Compare the same
initial population and arrivals over a fixed common horizon \(T\) that clears
all compared plans. Here \(J\) is total system time, not mean time.

\textbf{1. Start with conservation and complete time accounting.}
Internal transfers cancel in the population total, so
\[
 \dot P(t;\theta)=A(t)-E(t;\theta),\qquad
 J(\theta)=\int_0^T P(t;\theta)\dd t.
\]
Because initial mass and exogenous arrivals do not depend on the meter setting,
\[
 \frac{\partial P(t;\theta)}{\partial\theta}
 =-\int_0^t\frac{\partial E(u;\theta)}{\partial\theta}\dd u.
\]
Counting external waiting is essential: metering changes internal release,
not the exogenous arrival population \(A\).

\textbf{2. Differentiate the cost and reverse the integration order.}
An exit-rate change at time \(u\) changes remaining population at every later
time from \(u\) to \(T\), giving duration weight \(T-u\):
\begin{align*}
 \frac{\partial J}{\partial\theta}
 &=-\int_0^T\int_0^t
       \frac{\partial E(u;\theta)}{\partial\theta}\dd u\dd t\\
 &=-\int_0^T\int_u^T
       \frac{\partial E(u;\theta)}{\partial\theta}\dd t\dd u
 =\boxed{-\int_0^T(T-u)
       \frac{\partial E(u;\theta)}{\partial\theta}\dd u.}
\end{align*}
This is the continuous accounting formula on slide 13. For CTM timestep
\(h\), \(K\) steps and left-endpoint accounting, the exact discrete companion is
\[
 P^{k+1}=P^k+h(A^k-E^k),\qquad J_h=h\sum_{k=0}^{K-1}P^k,
 \qquad
 \frac{\partial J_h}{\partial\theta}
 =-h^2\sum_{k=0}^{K-2}(K-1-k)\frac{\partial E^k}{\partial\theta}.
\]

\textbf{3. Obtain exit sensitivity by the chain rule through the traffic model.}
The accounting identity does not itself determine how exits respond. In the
CTM, the meter schedule \(r^k(\theta)\) feeds release/storage, lateral exchange,
forward cell flows and merge allocation. Write the whole state update as
\(x^{k+1}=F(x^k,r^k)\) and final exit rate as \(E^k=H(x^k,r^k)\).
With \(z^k=\partial x^k/\partial\theta\) and fixed initial state,
\[
 z^0=0,\qquad
 z^{k+1}=F_xz^k+F_r\frac{\partial r^k}{\partial\theta},\qquad
 \frac{\partial E^k}{\partial\theta}
 =H_xz^k+H_r\frac{\partial r^k}{\partial\theta}.
\]
These sensitivities go into the weighted-exit formula. Moving final departures
earlier reduces total time. Moving waiting between ramp and mainline while
leaving total final exits unchanged gives zero total-time change.

\textbf{Limits.} Derivatives above assume an integrable sensitivity and fixed
smooth active branches; min/max ties require separate event treatment. A
blocked rate change can have zero sensitivity. Mean time divides both cost and
derivative by fixed total mass. This is an analytical explanation, not numerical
gradient verification or independent usefulness evidence.

\small Sources: current RM9 separate-resource/observer audit, sensitivity
section; RM8 task/gradient screen.
No model or review-gate decision is changed.
\end{document}
